\chapter{Introducción}

La publicación de \textit{A Mathematical Theory of Communication} por Claude E. Shannon en 1948 marcó el nacimiento definitivo de la teoría de la información y la era digital contemporánea \parencite{shannon1948mathematical}. Este trabajo pionero estableció los límites teóricos absolutos para la compresión de datos y su transmisión a través de canales con y sin ruido, sentando las bases matemáticas para las telecomunicaciones modernas.

En el presente desarrollo, se exploran y analizan los cuestionamientos centrales derivados de esta teoría fundamental. Entre los temas abordados se encuentran el intercambio entre ancho de banda y la relación señal/ruido en métodos de modulación por pulsos, la naturaleza inherentemente logarítmica de la medición de la información, el comportamiento probabilístico de la entropía en sistemas binarios, y las implicaciones que tiene la capacidad del canal sobre la velocidad y redundancia en la transmisión de mensajes.
